% Options for packages loaded elsewhere
% Options for packages loaded elsewhere
\PassOptionsToPackage{unicode}{hyperref}
\PassOptionsToPackage{hyphens}{url}
\PassOptionsToPackage{dvipsnames,svgnames,x11names}{xcolor}
%
\documentclass[
  german,
  10pt,
  twoside]{article}
\usepackage{xcolor}
\usepackage[top=2.5cm,bottom=2.5cm,left=2cm,right=2cm]{geometry}
\usepackage{amsmath,amssymb}
\setcounter{secnumdepth}{5}
\usepackage{iftex}
\ifPDFTeX
  \usepackage[T1]{fontenc}
  \usepackage[utf8]{inputenc}
  \usepackage{textcomp} % provide euro and other symbols
\else % if luatex or xetex
  \usepackage{unicode-math} % this also loads fontspec
  \defaultfontfeatures{Scale=MatchLowercase}
  \defaultfontfeatures[\rmfamily]{Ligatures=TeX,Scale=1}
\fi
\usepackage{lmodern}
\ifPDFTeX\else
  % xetex/luatex font selection
\fi
% Use upquote if available, for straight quotes in verbatim environments
\IfFileExists{upquote.sty}{\usepackage{upquote}}{}
\IfFileExists{microtype.sty}{% use microtype if available
  \usepackage[]{microtype}
  \UseMicrotypeSet[protrusion]{basicmath} % disable protrusion for tt fonts
}{}
\usepackage{setspace}
\makeatletter
\@ifundefined{KOMAClassName}{% if non-KOMA class
  \IfFileExists{parskip.sty}{%
    \usepackage{parskip}
  }{% else
    \setlength{\parindent}{0pt}
    \setlength{\parskip}{6pt plus 2pt minus 1pt}}
}{% if KOMA class
  \KOMAoptions{parskip=half}}
\makeatother
% Make \paragraph and \subparagraph free-standing
\makeatletter
\ifx\paragraph\undefined\else
  \let\oldparagraph\paragraph
  \renewcommand{\paragraph}{
    \@ifstar
      \xxxParagraphStar
      \xxxParagraphNoStar
  }
  \newcommand{\xxxParagraphStar}[1]{\oldparagraph*{#1}\mbox{}}
  \newcommand{\xxxParagraphNoStar}[1]{\oldparagraph{#1}\mbox{}}
\fi
\ifx\subparagraph\undefined\else
  \let\oldsubparagraph\subparagraph
  \renewcommand{\subparagraph}{
    \@ifstar
      \xxxSubParagraphStar
      \xxxSubParagraphNoStar
  }
  \newcommand{\xxxSubParagraphStar}[1]{\oldsubparagraph*{#1}\mbox{}}
  \newcommand{\xxxSubParagraphNoStar}[1]{\oldsubparagraph{#1}\mbox{}}
\fi
\makeatother


\usepackage{graphicx}
\makeatletter
\newsavebox\pandoc@box
\newcommand*\pandocbounded[1]{% scales image to fit in text height/width
  \sbox\pandoc@box{#1}%
  \Gscale@div\@tempa{\textheight}{\dimexpr\ht\pandoc@box+\dp\pandoc@box\relax}%
  \Gscale@div\@tempb{\linewidth}{\wd\pandoc@box}%
  \ifdim\@tempb\p@<\@tempa\p@\let\@tempa\@tempb\fi% select the smaller of both
  \ifdim\@tempa\p@<\p@\scalebox{\@tempa}{\usebox\pandoc@box}%
  \else\usebox{\pandoc@box}%
  \fi%
}
% Set default figure placement to htbp
\def\fps@figure{htbp}
\makeatother



\ifLuaTeX
\usepackage[bidi=basic]{babel}
\else
\usepackage[bidi=default]{babel}
\fi
% get rid of language-specific shorthands (see #6817):
\let\LanguageShortHands\languageshorthands
\def\languageshorthands#1{}
\ifLuaTeX
  \usepackage[german]{selnolig} % disable illegal ligatures
\fi


\setlength{\emergencystretch}{3em} % prevent overfull lines

\providecommand{\tightlist}{%
  \setlength{\itemsep}{0pt}\setlength{\parskip}{0pt}}



 
\usepackage[style=authoryear,backend=biber,style=authoryear-comp,uniquename=false,uniquelist=false]{biblatex}
\addbibresource{../references.bib}


% Font configuration (Statis Sans alternative - use sans-serif throughout)
\renewcommand{\familydefault}{\sfdefault}
\usepackage{helvet}  % Use Helvetica-like font

\usepackage{csquotes}
\usepackage{booktabs}
\usepackage{array}
\usepackage{multirow}
\usepackage{wrapfig}
\usepackage{float}
\usepackage{colortbl}
\usepackage{pdflscape}
\usepackage{threeparttable}
\usepackage{threeparttablex}
\usepackage[normalem]{ulem}
\usepackage{makecell}
\usepackage{xcolor}
\usepackage{multicol}
\usepackage{fancyhdr}
\usepackage{titlesec}
\usepackage{etoolbox}
\usepackage{tabularx}
\usepackage{afterpage}
\usepackage{tikz}
\usepackage{longtable}
\usepackage{totcount}
\usepackage{xstring}

% Time display with Berlin timezone (+1 hour offset)
\newcount\berlinhr
\newcount\berlinmin
\newcommand{\berlintime}{%
  \berlinhr=\time\relax
  \divide\berlinhr by 60\relax
  \berlinmin=\time\relax
  \multiply\berlinhr by -60\relax
  \advance\berlinmin by \berlinhr\relax
  \berlinhr=\time\relax
  \divide\berlinhr by 60\relax
  \advance\berlinhr by 1\relax
  \ifnum\berlinhr>23\relax
    \advance\berlinhr by -24\relax
  \fi
  \ifnum\berlinhr<10 0\fi\the\berlinhr:\ifnum\berlinmin<10 0\fi\the\berlinmin%
}

% Register counters for total counting
\regtotcounter{table}
\regtotcounter{figure}

% Suppress default maketitle
\renewcommand{\maketitle}{}

% ===== DRAFT MODE CONTROL =====
% Set \draftmodetrue to enable draft mode (metadata page + ENTWURF header)
% Set \draftmodefalse to disable draft mode (for final version)
\newif\ifdraftmode
\draftmodetrue  % <-- Change to \draftmodefalse to disable draft mode

% Define default commands for counts (will be overridden by filter)
\providecommand{\totalchars}{0}
\providecommand{\totalwords}{0}
\providecommand{\abstractchars}{0}
\providecommand{\abstractwords}{0}
\providecommand{\zusammenfassungchars}{0}
\providecommand{\zusammenfassungwords}{0}
\providecommand{\footnotecountnum}{0}
\providecommand{\citationcountnum}{0}
\providecommand{\schlusselwortercount}{0}
\providecommand{\keywordscount}{0}

% Define metadata commands from Pandoc variables
\newcommand{\doctitle}{$title$}
\newcommand{\docauthors}{$author-short$}
\newcommand{\docpublication}{$publication$}
\newcommand{\docpublicationissue}{$publication-issue$}
\newcommand{\docshorttitle}{$short-title$}

\ifdraftmode
% Define metadata page command
\newcommand{\metadatapage}{%
  \clearpage
  \thispagestyle{empty}
  \onecolumn
  \vspace*{1.5cm}
  
  % Header section
  {\raggedright
  {\LARGE\bfseries\color{wistaBlue} Dokument-Metadaten}\\[0.5cm]
  
  \begin{tabular}{@{}ll}
    \textbf{Titel:} & \doctitle \\[0.3em]
    \textbf{Autoren:} & \docauthors \\[0.3em]
    \textbf{Veröffentlichung:} & \docpublication\ \docpublicationissue \\[0.3em]
    \textbf{Erstellungsdatum:} & \today\ um \berlintime\ Uhr \\[0.3em]
  \end{tabular}
  
  \vspace{0.5cm}
  {\color{wistaBlue}\rule{\textwidth}{1pt}}
  \par}
  
  \vspace{1cm}
  
  % Dokumentstatistik
  {\large\bfseries\color{wistaDarkGray} Dokumentstatistik}\\[0.5cm]
  
  \small
  \begin{tabular}{@{}lrrrrrr@{}}
    \toprule
    \textbf{Abschnitt} & \textbf{Wörter} & \textbf{Zeichen} & \textbf{Fußn.} & \textbf{Quellen} & \textbf{Vorgabe} \\
    \midrule
    Insgesamt & \textcolor{wistaBlue}{\textbf{\totalwords}} & \textcolor{wistaBlue}{\textbf{\totalchars}} & \textcolor{wistaBlue}{\textbf{\footnotecountnum}} & \textcolor{wistaBlue}{\textbf{\citationcountnum}} & 30\,000 Z. \\
    Zusammenfassung & \textcolor{wistaBlue}{\textbf{\zusammenfassungwords}} & \textcolor{wistaBlue}{\textbf{\zusammenfassungchars}} & --- & --- & 800 Z. \\
    Abstract & \textcolor{wistaBlue}{\textbf{\abstractwords}} & \textcolor{wistaBlue}{\textbf{\abstractchars}} & --- & --- & 700 Z. \\
    \midrule
    \sectionrows
    \bottomrule
  \end{tabular}
  \normalsize
  
  \vspace{0.5cm}
  
  % Keywords list
  {\large\bfseries\color{wistaDarkGray} Schlüsselwörter}\\[0.3cm]
  \small
  \schlusselworterlist
  \normalsize
  
  \vspace{0.5cm}
  
  {\large\bfseries\color{wistaDarkGray} Keywords}\\[0.3cm]
  \small
  \keywordslist
  \normalsize
  
  \vspace{1cm}
  
  % Tables list
  {\large\bfseries\color{wistaDarkGray} Tabellen}\\[0.3cm]
  
  \small
  \begin{tabular}{@{}lp{12cm}@{}}
    \toprule
    \textbf{Nr.} & \textbf{Beschreibung} \\
    \midrule
    \tablerows
    \bottomrule
  \end{tabular}
  \normalsize
  
  \vspace{1cm}
  
  % Figures list
  {\large\bfseries\color{wistaDarkGray} Abbildungen}\\[0.3cm]
  
  \small
  \begin{tabular}{@{}lp{12cm}@{}}
    \toprule
    \textbf{Nr.} & \textbf{Beschreibung} \\
    \midrule
    \figurerows
    \bottomrule
  \end{tabular}
  \normalsize
  
  \clearpage
}
\fi

% Define WISTA colors (from Kristiansen Word template)
\definecolor{wistaBlue}{RGB}{91,155,213}
\definecolor{wistaLightBlue}{RGB}{120,170,215}
\definecolor{wistaDarkGray}{RGB}{96,94,92}
\definecolor{wistaLightGray}{RGB}{128,128,128}

% Section formatting (##) - number 20% bigger, lines closer, text closer to bottom line
% e.g., "2" with lines around "Einleitung"
\titleformat{\section}[display]
  {\normalfont\large\bfseries}
  {\color{wistaBlue}\Large\bfseries\thesection\\\noindent\rule{\columnwidth}{1.5pt}}
  {0pt}
  {\vspace{0pt}\color{wistaDarkGray}\large\bfseries}
  [\vspace{0pt}\noindent{\color{wistaBlue}\rule{\columnwidth}{1.5pt}}\vspace{0.05cm}]
\titlespacing*{\section}{0pt}{1.5ex plus 0.3ex minus .1ex}{0.2ex plus .1ex}

% Subsection formatting (###) - grey text with blue line under
% e.g., "2.1 Handlungsbedarf"
\renewcommand{\thesubsection}{\thesection.\arabic{subsection}}
\titleformat{\subsection}
  {\color{wistaDarkGray}\normalsize\bfseries}
  {\color{wistaDarkGray}\normalsize\bfseries\thesubsection}
  {0.5em}
  {\color{wistaDarkGray}\normalsize\bfseries}
  [\vspace{0pt}\noindent{\color{wistaBlue}\rule{\columnwidth}{0.8pt}}\vspace{0.05cm}]
\titlespacing*{\subsection}{0pt}{1ex plus 0.2ex minus .1ex}{0.2ex plus .1ex}
  
% Subsubsection formatting (####) - grey text, no underline
% e.g., "2.1.1 Steigende Datenmenge"
\renewcommand{\thesubsubsection}{\thesubsection.\arabic{subsubsection}}
\titleformat{\subsubsection}
  {\color{wistaDarkGray}\normalsize\bfseries}
  {\color{wistaDarkGray}\normalsize\bfseries\thesubsubsection}
  {0.5em}
  {\color{wistaDarkGray}\normalsize\bfseries}

% Paragraph formatting (4th level) - numbered as subsection.subsubsection.paragraph
\renewcommand{\theparagraph}{\thesubsubsection.\arabic{paragraph}}
\titleformat{\paragraph}[runin]
  {\color{wistaDarkGray}\normalsize\bfseries}
  {\color{wistaBlue}\normalsize\bfseries\theparagraph}
  {0.5em}
  {\color{wistaDarkGray}\normalsize\bfseries}
  [.]
\titlespacing*{\paragraph}{0pt}{1ex plus 0.5ex minus 0.2ex}{0.5em}

% Header and footer setup with blue line under header
\pagestyle{fancy}
\fancyhf{}
\fancyhead[L]{\small\textit{\docauthors}}
\ifdraftmode
  \fancyhead[C]{\small\color{red}\textbf{ENTWURF} -- \today}
\fi
\fancyhead[R]{\small\textit{\docshorttitle}}
\fancyfoot[L]{\small\thepage}
\fancyfoot[R]{\small Statistisches Bundesamt | \docpublication\ | {\color{wistaBlue}\docpublicationissue}}
\renewcommand{\headrulewidth}{0pt}
\renewcommand{\footrulewidth}{0pt}
% Add blue line under header
\renewcommand{\headrule}{\color{wistaBlue}\hrule width\headwidth height 1pt \vskip 2pt}
\renewcommand{\headrulewidth}{1pt}

% Make all figures span both columns (full width)
\let\origfigure\figure
\let\endorigfigure\endfigure
\renewenvironment{figure}[1][htbp]{%
  \origfigure*[#1]%
}{%
  \endorigfigure*%
}

% Allow manual control of table* environment for full-width tables
% Tables by default stay in single column, but can use table* manually



% Make bibliography full-width (one column) - for biblatex
\defbibheading{bibliography}[\bibname]{%
  \onecolumn
  \section*{#1}%
  \markboth{#1}{#1}%
}



% Footnotes stay within column
\usepackage[stable]{footmisc}

% Custom footnote formatting: "|^1" in dark WISTA gray
\renewcommand{\thefootnote}{{\color{wistaDarkGray}|$^{\arabic{footnote}}$}}

% Pifont for dingbat symbols
\usepackage{pifont}

% Better table handling for two-column mode
\usepackage{array,tabularx,calc}

% Custom WISTA table format
\usepackage{caption}
\usepackage{xcolor}
\usepackage{colortbl}
\usepackage{array}

% Define custom caption format for WISTA tables
\DeclareCaptionFormat{wista}{%
  {\color{wistaLightBlue}\bfseries #1#2}\par\vspace{0.2em}%
  {\color{wistaDarkGray}\normalfont #3}\par\vspace{0.3em}%
}

% Set up table captions with custom WISTA format
\captionsetup[table]{
  position=top,
  format=wista,
  justification=justified,
  singlelinecheck=false,
  labelsep=none,
  skip=0.3em
}

% Better table spacing and white inner lines
\setlength{\tabcolsep}{6pt}
\renewcommand{\arraystretch}{1.15}

% Set white lines for tables by default, blue for first column
\AtBeginDocument{
  \arrayrulecolor{white}
  \setlength{\arrayrulewidth}{0.3pt}
}

% Define commands for colored table borders
\newcommand{\bluevrule}{\color{wistaBlue}\vrule}
\newcommand{\whitevrule}{\color{white}\vrule}

% Define WISTA table colors
\definecolor{wistaTableBg}{RGB}{240,248,255}     % Very light blue background
\definecolor{wistaHeaderBg}{RGB}{220,230,240}    % Light blue header background
\definecolor{wistaFirstCol}{RGB}{65,125,180}     % Blue for first column lines

% Custom WISTA table command
\newcommand{\wistaTableStart}{%
  \rowcolors{2}{wistaTableBg}{wistaTableBg}%
  \arrayrulecolor{wistaFirstCol}%
}

% Custom commands for table styling
\newcommand{\wistaHeaderRow}[1]{%
  \rowcolor{wistaHeaderBg}%
  \arrayrulecolor{wistaDarkGray}%
  #1 \\
  \arrayrulecolor{white}%
}

\newcommand{\wistaFirstColRule}{\arrayrulecolor{wistaFirstCol}}
\newcommand{\wistaWhiteRule}{\arrayrulecolor{white}}

% Custom bullet points with > character - 35% width, 95% height, black extra bold, half line spacing
\usepackage{enumitem}
\setlist[itemize]{label={\scalebox{0.35}[0.95]{\fontseries{bx}\selectfont >}}, leftmargin=1.5em, itemsep=5pt, topsep=3pt, parsep=0pt}

% Redefine ref to include pifont 247 (arrow) for figure/table/section references
\let\oldref\ref
\renewcommand{\ref}[1]{\oldref{#1}\raisebox{-0.2ex}{\scriptsize\color{wistaBlue}\ding{247}}}
\makeatletter
\@ifpackageloaded{caption}{}{\usepackage{caption}}
\AtBeginDocument{%
\ifdefined\contentsname
  \renewcommand*\contentsname{Inhaltsverzeichnis}
\else
  \newcommand\contentsname{Inhaltsverzeichnis}
\fi
\ifdefined\listfigurename
  \renewcommand*\listfigurename{Abbildungsverzeichnis}
\else
  \newcommand\listfigurename{Abbildungsverzeichnis}
\fi
\ifdefined\listtablename
  \renewcommand*\listtablename{Tabellenverzeichnis}
\else
  \newcommand\listtablename{Tabellenverzeichnis}
\fi
\ifdefined\figurename
  \renewcommand*\figurename{Abbildung}
\else
  \newcommand\figurename{Abbildung}
\fi
\ifdefined\tablename
  \renewcommand*\tablename{Tabelle}
\else
  \newcommand\tablename{Tabelle}
\fi
}
\@ifpackageloaded{float}{}{\usepackage{float}}
\floatstyle{ruled}
\@ifundefined{c@chapter}{\newfloat{codelisting}{h}{lop}}{\newfloat{codelisting}{h}{lop}[chapter]}
\floatname{codelisting}{Listing}
\newcommand*\listoflistings{\listof{codelisting}{Listingverzeichnis}}
\makeatother
\makeatletter
\makeatother
\makeatletter
\@ifpackageloaded{caption}{}{\usepackage{caption}}
\@ifpackageloaded{subcaption}{}{\usepackage{subcaption}}
\makeatother
\usepackage{bookmark}
\IfFileExists{xurl.sty}{\usepackage{xurl}}{} % add URL line breaks if available
\urlstyle{same}
\hypersetup{
  pdftitle={EFFIZIENTE AUSWERTUNG DES BUSINESS-TAX-PANELS MITHILFE VON R},
  pdfauthor={Oliver Hauke; Annette Kristiansen},
  pdflang={de},
  colorlinks=true,
  linkcolor={blue},
  filecolor={Maroon},
  citecolor={Blue},
  urlcolor={Blue},
  pdfcreator={LaTeX via pandoc}}


\title{EFFIZIENTE AUSWERTUNG DES BUSINESS-TAX-PANELS MITHILFE VON R}
\author{Oliver Hauke \and Annette Kristiansen}
\date{2025-12-05}
\begin{document}
\maketitle

% Custom WISTA-style title page
\begin{titlepage}
\thispagestyle{empty}
\pagestyle{empty}
\onecolumn

% Horizontal blue line at top (header line)
\noindent{\color{wistaBlue}\rule{\textwidth}{1pt}}

\vspace{-0.5cm}

% Vertical blue line on the left (6.3cm from left edge, starts below header line, stops above footer)
\begin{tikzpicture}[remember picture, overlay]
  \draw[wistaBlue, line width=2pt] 
    ([xshift=6.3cm, yshift=-4cm]current page.north west) -- 
    ([xshift=6.3cm, yshift=3.5cm]current page.south west);
\end{tikzpicture}

\vspace*{0.5cm}

% Left side - short author bios (right-aligned to the vertical line)
\begin{minipage}[t]{3.9cm}
\selectlanguage{ngerman}%
\raggedleft
\hyphenpenalty=0
\tolerance=1000
\fontsize{7}{8.5}\selectfont
\vspace{0pt}%
{\fontsize{8}{9.5}\selectfont\color{wistaBlue}\textbf{Oliver Hauke}}\\[0.15cm]
ist Mathematiker und IT-Referent im Kompetenzzentrum „Analyse und Auswertung" des Statistischen Bundesamtes. Er verantwortet die Weiterentwicklung der technischen Infrastruktur des Forschungsdatenzentrums und berät das Netzwerk empirische Steuerforschung in der effizienten Nutzung der Systeme.\\[0.5cm]

{\fontsize{8}{9.5}\selectfont\color{wistaBlue}\textbf{Annette Kristiansen}}\\[0.15cm]
ist Ökonomin und Referentin im Referat „Unternehmenssteuern" des Statistischen Bundesamtes. Sie beschäftigt sich mit der Zusammenführung der Unternehmenssteuerstatistiken und betreut das Netzwerk empirische Steuerforschung. Für ihre externe Promotion an der Universität Bremen untersucht sie die technologische Vielfalt multinationaler Unternehmen.

\end{minipage}%
\hspace{0.4cm}% Gap to vertical line (left side)
\hspace{0.4cm}% Gap from vertical line (right side)
% Right side - title, authors, keywords, abstracts
\begin{minipage}[t]{11.8cm}
\raggedright
\vspace{1.2cm}%

% Title (uppercase via command, not bold)
{\LARGE\color{wistaDarkGray}\begin{spacing}{1.3}\MakeUppercase{Effiziente Auswertung des Business-Tax-Panels mithilfe von R}\end{spacing}}

\vspace{-0.2cm}

% Blue horizontal line (underline for title)
{\color{wistaBlue}\rule{\linewidth}{0.5pt}}

\vspace{0.5cm}

% Authors inline, dark gray
{\large\color{wistaDarkGray} Oliver Hauke, Annette Kristiansen}

\vspace{0.05cm}

% Blue line (underline for authors)
{\color{wistaBlue}\rule{\linewidth}{0.5pt}}

\vspace{0.5cm}

% Keywords (pifont 247 arrow then bold label and blue text)
\hangindent=1.2em\hangafter=1
{\color{wistaBlue}\ding{247}\ {\bfseries Schlüsselwörter:} Effiziente Analysen, R-Programmierung, Forschungsdaten, Unternehmenssteuerstatistik, Business-Tax-Panel, Netzwerk empirische Steuerforschung}

\vspace{0.3cm}

% Zusammenfassung (uppercase via command)
{\color{wistaDarkGray}\bfseries\MakeUppercase{Zusammenfassung}}\\[0.15cm]
\small
Das Forschungsdatenzentrum des Statistischen Bundesamtes baut seine technische Infrastruktur gezielt aus, um der Wissenschaft erweiterte Analysemöglichkeiten amtlicher Daten bereitzustellen. Seit November 2025 steht Forschenden exklusiver Zugang zu diesen erweiterten Systemressourcen in R und Stata zur Verfügung. R-basierte Tools -- insbesondere Apache Arrow, Parquet und DuckDB -- ermöglichen es, gängige Analysen großer Forschungsdatensätze in Echtzeit auszuführen. Am Beispiel des Business-Tax-Panel (2013 bis 2019) werden die erzielbaren Effizienzgewinne bei der Auswertung umfangreicher steuerstatistischer Daten im Rahmen des Netzwerks empirische Steuerforschung demonstriert.

\vspace{0.5cm}

% English keywords and abstract (italicized)
\hangindent=1.2em\hangafter=1
{\itshape\color{wistaBlue}\ding{247}\ {\bfseries Keywords:} Efficient analysis -- R programming -- research data -- corporate tax statistics -- Business Tax Panel -- empirical tax research network}

\vspace{0.3cm}

{\itshape\color{wistaDarkGray}\bfseries\MakeUppercase{Abstract}}\\[0.15cm]
\small\itshape
The Research Data Centre at the Federal Statistical Office is strategically expanding its technical infrastructure to provide researchers with enhanced analytical capabilities for official data. Since November 2025, researchers have had exclusive access to these extended system resources in R and Stata. R-based tools -- particularly Apache Arrow, Parquet, and DuckDB -- enable real-time analysis of large research datasets. Using the Business Tax Panel (2013--2019), this paper demonstrates the achievable efficiency gains in analysing large-scale tax statistics within the empirical tax research network.

\end{minipage}

\vfill

\vspace{0.5cm}
% Footer matching document style
\noindent\makebox[\textwidth]{%
  \small 1%
  \hfill%
  \small Statistisches Bundesamt | WISTA | {\color{wistaBlue}01/2026}%
}

\end{titlepage}

% Stay in one-column mode for TOC (will be switched later)
\pagestyle{fancy}

\renewcommand*\contentsname{Inhaltsverzeichnis}
{
\hypersetup{linkcolor=}
\setcounter{tocdepth}{3}
\tableofcontents
}

\setstretch{1}
\renewcommand{\totalchars}{4 521}
\renewcommand{\totalwords}{552}
\renewcommand{\abstractchars}{0}
\renewcommand{\abstractwords}{0}
\renewcommand{\zusammenfassungchars}{0}
\renewcommand{\zusammenfassungwords}{0}
\renewcommand{\footnotecountnum}{0}
\renewcommand{\citationcountnum}{0}
\renewcommand{\schlusselwortercount}{0}
\renewcommand{\keywordscount}{0}
\renewcommand{\doctitle}{EFFIZIENTE AUSWERTUNG DES BUSINESS-TAX-PANELS MITHILFE VON R}
\renewcommand{\docauthors}{Oliver Hauke, Annette Kristiansen}
\renewcommand{\docpublication}{WISTA}
\renewcommand{\docpublicationissue}{01/2026}
\renewcommand{\docshorttitle}{Effiziente Analyse großer Forschungsdaten}
\newcommand{\sectionrows}{            Einleitung & \textcolor{wistaBlue}{\textbf{1 231}} & \textcolor{wistaBlue}{\textbf{1 464}} & \textcolor{wistaBlue}{\textbf{0}} & \textcolor{wistaBlue}{\textbf{0}} & --- \\
            Anwendungsfall Business-Tax-Panel & \textcolor{wistaBlue}{\textbf{125}} & \textcolor{wistaBlue}{\textbf{148}} & \textcolor{wistaBlue}{\textbf{0}} & \textcolor{wistaBlue}{\textbf{0}} & --- \\
            Techniche Infrastruktur & \textcolor{wistaBlue}{\textbf{145}} & \textcolor{wistaBlue}{\textbf{172}} & \textcolor{wistaBlue}{\textbf{0}} & \textcolor{wistaBlue}{\textbf{0}} & --- \\
            Datenverarbeitungsmethoden & \textcolor{wistaBlue}{\textbf{810}} & \textcolor{wistaBlue}{\textbf{963}} & \textcolor{wistaBlue}{\textbf{0}} & \textcolor{wistaBlue}{\textbf{0}} & --- \\
            Performanzvergleich & \textcolor{wistaBlue}{\textbf{999}} & \textcolor{wistaBlue}{\textbf{1 188}} & \textcolor{wistaBlue}{\textbf{0}} & \textcolor{wistaBlue}{\textbf{0}} & --- \\
            Ergebnis für das vollständige BTP & \textcolor{wistaBlue}{\textbf{84}} & \textcolor{wistaBlue}{\textbf{99}} & \textcolor{wistaBlue}{\textbf{0}} & \textcolor{wistaBlue}{\textbf{0}} & --- \\
            Fazit & \textcolor{wistaBlue}{\textbf{407}} & \textcolor{wistaBlue}{\textbf{484}} & \textcolor{wistaBlue}{\textbf{0}} & \textcolor{wistaBlue}{\textbf{0}} & --- \\
}
\newcommand{\tablerows}{            \hyperref[tbl-pkg-versions]{Tabelle~1} & Versionen verfügbarer R-Packages \\
            \hyperref[tbl-baseline]{Tabelle~2} & Baselines in SAS und Stata (anhand 1\% des simuliertes BTP) \\
            \hyperref[tbl-benchmark-small]{Tabelle~3} & Performanz Messungen für die Fallzahlberechnung anhand 1\% des BTP \\
            \hyperref[tbl-benchmark-regr-small]{Tabelle~4} & Performanz Messungen für die Regressionsberechnung anhand 1\% des BTP \\
            \hyperref[tbl-benchmark]{Tabelle~5} & Performanz Messungen für die Fallzahlberechnung anhand 10\% des BTP \\
            \hyperref[tbl-benchmark-regr]{Tabelle~6} & Performanz Messungen für die Regressionsberechnung anhand 10\% des BTP \\
            \hyperref[tbl-bm-cnt-small]{Tabelle~7} & Benchmarkergebnisse für Fallzahlen (kleine Stichproben des BTP) \\
            \hyperref[tbl-bm-reg-small]{Tabelle~8} & Benchmarkergebnisse für Regression (kleine Stichproben des BTP) \\
}
\newcommand{\figurerows}{            \hyperref[fig-optimization-strategies]{Abbildung~1} & Strategien zur Performanz-Optimierung im FDZ-Bund \\
}
\newcommand{\schlusselworterlist}{}
\newcommand{\keywordslist}{}
\ifdraftmode
\metadatapage
\fi

\setcounter{subsection}{5}

\section{Fazit}\label{sec-ausblick}

\subsection{Kernaussage des Artikels}\label{kernaussage-des-artikels}

Wir empfehlen eine Datenhaltung im Parquet-Format und \textbf{R}
incl.~\{arrow\} + \{parquet\} (bei Bedarf \{duckdb\}, \{polars\}) zu
erproben und zu etablieren! Denn unsere Untersuchung demonstriert ganz
neue Möglichkeiten, analytische Fragestellungen anhand großer
Datenmengen in Echtzeit zu beantworten.

\subsection{Auswirkung auf IT-Infrastruktur und
Beratung}\label{auswirkung-auf-it-infrastruktur-und-beratung}

\begin{itemize}
\tightlist
\item
  Fokus auf Hardware und Software mit größtem Mehrwert (weniger RAM und
  R anstatt Spark, GPUs, Datenbanken, Cloud, \ldots)
\item
  Fokus auf Maintainance und Patching von R, Parquet und hier genannte
  Packages \{arrow\}, \{duckdb\}, \{polars\}
\item
  Etablierung dieser Softwarelösungen kann langfristig\ldots{}

  \begin{itemize}
  \tightlist
  \item
    die Hardwareanforderungen reduzieren
  \item
    die Analysesoftwarevielfalt verschlanken, somit Beschaffungen
    einsparen (besserer Ausdruck?)
  \end{itemize}
\item
  Schulungsangebot benötigt für \emph{``hoch-performante Auswertungen in
  R''}
\item
  eigener Beitrag zu Open Source entscheidend, um die Etablierung zu
  fördern.
\end{itemize}

\subsection{Auswirkungen auf die
Statistikproduktion}\label{auswirkungen-auf-die-statistikproduktion}

\begin{itemize}
\tightlist
\item
  Steigerung der Effizienz und Reaktionsgeschwindigkeit
\item
  (Vorsichtig!? Annette, lieber weg oder rein?) Vereinfachung der
  Statistikproduktion, durch leicht nachvollziehbaren Code und Fokus auf
  Wissensaufbau in R
\item
  Aufbau eines Datenbestands in Parquet
\item
  Anwendung auf weitere Produkte

  \begin{itemize}
  \tightlist
  \item
    Überarbeitung des BTP
  \item
    Mindeststeuer (oder?)
  \item
    DRG
  \item
    TPP 2
  \end{itemize}
\end{itemize}

\subsection{Auswirkungen auf die FDZ}\label{auswirkungen-auf-die-fdz}

\begin{itemize}
\tightlist
\item
  Berücksichtigung in Mustersyntaxen und Nutzendenberatung im FDZ Bund.
  Teilen von Code-Snippets und Templates fördern.
\item
  mögliche Vereinfachung der Betreuung (Fokus auf R, Reduktion von
  Performanz- und Programmierproblemen)
\item
  Bereitstellung von Parquet in geeigneter Aufteilung empfohlen
\item
  Erprobung im FDZ Bund Vorbild für andere FDZ
\item
  Optimierung von Hardware und Software Angebot
\end{itemize}

\subsection{Auswirkung auf die
Wissenschaft}\label{auswirkung-auf-die-wissenschaft}

\begin{itemize}
\tightlist
\item
  Verwendung von R bietet am FDZ Bund einen deutlichen
  \textbf{Wettbewerbsvorteil in der emp. Forschung}. Lediglich R ist (am
  GWAP und via KDFV) am FDZ Bund in der Lage große Datenmengen innerhalb
  von Sekunden zu analysieren und Datenmengen überhalb des verfügbaren
  Arbeitspeicher, wie das BTP, ohne vorherige Variablenselektion
  auszuwerten.

  \begin{itemize}
  \tightlist
  \item
    Vollmaterial kann statt Stichproben angeboten werden. Erleichtert
    die FDZ Betreuung und GWAP/KDFV-Verfügbarkeit
  \item
    mehr explorative oder komplexe Analysen sind in kürzerer Zeit
    möglich.
  \end{itemize}
\item
  Dazu ist R als frei verfügbare Open-Source Software leicht zu lernen
  und \{arrow\} bietet hoch-performante Analysen selbst für
  R-Basiskenntnisse. Empfehlung an die emp. Wissenschaft, sich verstärkt
  mit R zu befassen.
\item
  Für analoge über-RAM-Verarbeitungen bietet Stata aktuell keine Lösung.
  SAS bietet bei weitem nicht hier gemessene Performanz von R. R
  hervorragende Möglichkeiten bietet Performanz und
  Anwenderfreundlichkeit zu übertreffen.
\item
  Weiteres Vorgehen: \textbf{NeSt-Pilotprojekt}. Aus dem realen Use Case
  Optimierung der Systemkonfiguration und Beratungsbedarfe
  identifizieren.
\end{itemize}

\subsection{Schlussbemerkung}\label{schlussbemerkung}

Die Analyse großer Datensätze in der amtlichen Statstik und den
Forschungsdatenzen können von einem signifikanten Performanz-Boost
profitieren, wenn die vorgestellten Werkzeuge adaptiert werden.
\textbf{Arrow, DuckDB, Polars und Parquet} -- bieten das Potenzial, die
Effizienz erheblich zu steigern und neue Analysemöglichkeiten zu
eröffnen.

Der Übergang erfordert jedoch eine \textbf{koordinierte Anstrengung} von
Infrastrukturbetreibern, amtlicher Statistik, Forschungsdatenzentren und
Forschenden. Die systematische Wissensaufbau, Dokumentation und
Verbreitung von Best Practices wird entscheidend für den Erfolg sein.

Die amtliche Statistik kann von diesen Entwicklungen in mehrfacher
Hinsicht profitieren: \textbf{schnellere Erkenntnisgewinnung},
\textbf{umfassendere Analysen} und \textbf{effizientere
Ressourcennutzung}. Dies trägt letztlich zur Stärkung der
evidenzbasierten Politik- und Gesellschaftsberatung bei.


\printbibliography



\end{document}
